\section{Conclusion}
\label{sec:conclusion}

This study compared TimesNet and DilatedRNN under common, leakage-controlled
hourly, daily, and monthly GHI tasks at ten electric-vehicle manufacturing
locations. Inputs, partitions, origins, post-processing, metrics, and location
weights were matched within each task. No architecture was universally
superior.

TimesNet had lower macro-MAE at 24 and 72 hours and through 14 days;
DilatedRNN led at 48 hours, 30 days, the primary one-month task, and the
three- and six-month exploratory prefixes. Terminal wins favored DilatedRNN at
six of ten hourly, daily, and one-month sites and seven of ten six-month sites,
although a few larger TimesNet gains sometimes reversed the macro average.
XGBoost led the full hourly and daily benchmarks, whereas calendar climatology
led the monthly tasks, so the neural comparison does not imply benchmark-wide
dominance.

Daily and monthly five-seed results support these patterns, but several
architecture differences were small relative to seed variability. The hourly
extension has only seed 42, multi-step targets overlap, and parameter budgets
differ. Evidence therefore supports resolution- and horizon-specific behaviour
rather than an intrinsic architectural ranking. Opposite weather contexts for
the strongest TimesNet gain and deficit likewise preclude a causal reading of
post-hoc meteorology.

Future work should cover prospective years and unseen sites, repeat the hourly
extension across seeds, report daylight-only metrics and measured cost, and
test causal weather inputs against the strong tabular and climatological
references. Released contracts, origin-level forecasts, seed metrics, and
integrity manifests support those extensions.
